\chapter*{Appendix B: Lagrangian and Hamiltonian Formalism for the RSVP Plenum}
\addcontentsline{toc}{chapter}{Appendix B: Lagrangian and Hamiltonian Formalism for the RSVP Plenum}

\section{Field Content and Kinematics}

Let $(\mathcal{M},g_{ab})$ be a four–dimensional, time–oriented Lorentzian manifold with signature
\[
(-,+,+,+),
\]
interpreted as the plenum of events. The RSVP theory is formulated in terms of three primary fields:
\begin{itemize}
  \item a scalar potential $\Phi : \mathcal{M}\to\mathbb{R}$,
  \item a vector field $v^a$ on $\mathcal{M}$ (the flow field),
  \item an entropy scalar $S : \mathcal{M}\to\mathbb{R}$.
\end{itemize}

Define the entropy current
\[
s^a := S v^a,
\]
and write $\nabla_a$ for the Levi–Civita connection of $g_{ab}$. The basic kinematic quantities are
\[
\theta := \nabla_a v^a,\qquad
\omega_{ab} := \nabla_{[a} v_{b]},\qquad
a^a := v^b \nabla_b v^a,
\]
together with the gradients $\nabla_a \Phi$ and $\nabla_a S$.

The curvature of the plenum is encoded in the Ricci scalar $R[g]$ and the Einstein tensor
\[
G_{ab} := R_{ab} - \tfrac{1}{2} R g_{ab}.
\]

\section{RSVP Lagrangian Density}

The total RSVP action is taken to be
\begin{equation}
\label{eq:RSVP-action}
\mathcal{S}_{\mathrm{RSVP}}[g,\Phi,v,S]
  = \int_{\mathcal{M}} \mathcal{L}_{\mathrm{RSVP}} \, d^4x
  = \int_{\mathcal{M}} \sqrt{-g}\,\Bigl(
      \mathcal{L}_{\mathrm{geom}}
    + \mathcal{L}_{\mathrm{fields}}
    + \mathcal{L}_{\mathrm{coupling}}
    \Bigr)\, d^4x.
\end{equation}

\subsection{Geometric term}

The geometric part is chosen as
\begin{equation}
\mathcal{L}_{\mathrm{geom}}
  = \frac{1}{2\kappa}\,R[g]
    - \Lambda_{\mathrm{eff}},
\end{equation}
where $\kappa>0$ is the gravitational coupling and $\Lambda_{\mathrm{eff}}$ is an effective cosmological term determined by the RSVP ground state. In the non–expanding plenum picture the global solution for $g_{ab}$ is stationary rather than FRW–expanding, but the local field equations follow the same variational structure.

\subsection{Field term}

The field Lagrangian collects kinetic and potential terms for $(\Phi,v^a,S)$:
\begin{equation}
\label{eq:L-fields}
\mathcal{L}_{\mathrm{fields}}
  = -\frac{\alpha_\Phi}{2}\,\nabla_a\Phi \nabla^a\Phi
    -\frac{\alpha_S}{2}\,\nabla_a S \nabla^a S
    -\frac{\alpha_v}{2}\,h_{ab} v^a v^b
    -\frac{\beta_v}{2}\,\omega_{ab}\omega^{ab}
    -V(\Phi,S),
\end{equation}
where:
\begin{itemize}
  \item $h_{ab} := g_{ab} + u_a u_b$ is the projector orthogonal to a chosen timelike unit vector $u^a$ (e.g.\ the preferred RSVP frame; in many applications one takes $u^a \propto v^a$ when $v^a$ is timelike),
  \item $\alpha_\Phi,\alpha_S,\alpha_v,\beta_v>0$ are coupling constants,
  \item $V(\Phi,S)$ is a scalar potential encoding the preferred ground state $(\Phi_\ast,S_\ast)$.
\end{itemize}
The sign choices ensure that for the Lorentzian signature $(-,+,+,+)$ the kinetic terms yield positive energy density in the Hamiltonian formalism.

\subsection{Coupling and entropy production}

The essential RSVP structure is encoded by coupling terms relating the divergence of $v^a$, the entropy gradient, and the scalar potential. A minimal choice is
\begin{equation}
\label{eq:L-coupling}
\mathcal{L}_{\mathrm{coupling}}
  =
  \lambda_1 \,\Phi\,\nabla_a v^a
  + \lambda_2 \,v^a\nabla_a S
  + \lambda_3 \,S\,\nabla_a v^a
  + \lambda_4 \,(\nabla_a S)(\nabla^a \Phi),
\end{equation}
with real couplings $\lambda_i$.

The combination of \eqref{eq:L-fields} and \eqref{eq:L-coupling} enforces, via the Euler–Lagrange equations, the characteristic RSVP relations:
\begin{itemize}
  \item divergence of $v^a$ as an entropic source,
  \item alignment between $\nabla_a S$ and the effective gravitational acceleration,
  \item monotonic entropy production along the flow $v^a$.
\end{itemize}

\section{Euler–Lagrange Equations}

Variation with respect to $\Phi$ gives
\begin{equation}
\label{eq:EOM-Phi}
\alpha_\Phi \nabla_a\nabla^a\Phi
  - \frac{\partial V}{\partial \Phi}
  + \lambda_1\,\nabla_a v^a
  + \lambda_4 \nabla_a\nabla^a S
  = 0.
\end{equation}

Variation with respect to $S$ yields
\begin{equation}
\label{eq:EOM-S}
\alpha_S \nabla_a\nabla^a S
  - \frac{\partial V}{\partial S}
  + \lambda_2 \nabla_a v^a
  + \lambda_2 v^a\nabla_a (\log\sqrt{-g})
  + \lambda_3 \nabla_a v^a
  + \lambda_4 \nabla_a\nabla^a \Phi
  = \sigma_{\mathrm{RSVP}},
\end{equation}
where $\sigma_{\mathrm{RSVP}}\ge 0$ is the entropy production density, introduced either as a Lagrange multiplier enforcing the second law, or as an emergent quantity derived from coarse–graining.

Variation with respect to $v^a$ gives
\begin{equation}
\label{eq:EOM-v}
\alpha_v h_{ab} v^b
  + \beta_v \nabla^b\omega_{ab}
  - \lambda_1\nabla_a\Phi
  - \lambda_2 \nabla_a S
  - \lambda_3 S a_a
  = 0.
\end{equation}
This equation expresses the balance between flow inertia, vorticity, and entropic driving forces.

Finally, variation with respect to $g_{ab}$ produces an effective Einstein equation
\begin{equation}
\label{eq:EOM-g}
\frac{1}{\kappa} G_{ab}
  + \Lambda_{\mathrm{eff}} g_{ab}
  = T^{\mathrm{RSVP}}_{ab},
\end{equation}
where the RSVP stress–energy tensor is
\begin{equation}
\begin{aligned}
T^{\mathrm{RSVP}}_{ab}
  &= \alpha_\Phi\left(\nabla_a\Phi\nabla_b\Phi
      - \tfrac{1}{2} g_{ab}\nabla_c\Phi\nabla^c\Phi\right)
   + \alpha_S\left(\nabla_a S\nabla_b S
      - \tfrac{1}{2} g_{ab}\nabla_c S\nabla^c S\right) \\
  &\quad
   + \alpha_v h_{c(a} h_{b)d} v^c v^d
   + \beta_v\left(\omega_{ac}\omega_b{}^c
      -\tfrac{1}{4} g_{ab} \omega_{cd}\omega^{cd}\right)
   + g_{ab} V(\Phi,S)
   + T^{\mathrm{coup}}_{ab},
\end{aligned}
\end{equation}
with $T^{\mathrm{coup}}_{ab}$ obtained from \eqref{eq:L-coupling}.  

Equations \eqref{eq:EOM-Phi}–\eqref{eq:EOM-g} define the RSVP dynamics at the Lagrangian level.

\section{3+1 Decomposition and Hamiltonian Density}

Choose a foliation of $\mathcal{M}$ by spacelike hypersurfaces $\Sigma_t$ with unit normal $n^a$, lapse $N$ and shift $N^i$. The metric decomposes as
\[
g_{ab} = -n_a n_b + h_{ab},\qquad
n_a n^a = -1,\qquad
h_{ab} n^b = 0,
\]
with $h_{ab}$ the spatial metric on $\Sigma_t$ and indices $i,j,k$ denoting spatial components.

Write the fields as
\[
\Phi(t,x^i),\quad S(t,x^i),\quad v^a = v_\perp n^a + v^i e_i{}^a,
\]
where $e_i{}^a$ is a spatial triad and $v_\perp := -n_a v^a$.

\subsection{Canonical momenta}

Define the canonical momenta conjugate to $\Phi$ and $S$:
\begin{align}
\Pi_\Phi
  &:= \frac{\partial\mathcal{L}_{\mathrm{RSVP}}}{\partial(\partial_t \Phi)}
   = \sqrt{h}\,\Bigl(
        \alpha_\Phi\,n^a \nabla_a\Phi
      + \lambda_4\,n^a \nabla_a S
     \Bigr), \\
\Pi_S
  &:= \frac{\partial\mathcal{L}_{\mathrm{RSVP}}}{\partial(\partial_t S)}
   = \sqrt{h}\,\Bigl(
        \alpha_S\,n^a \nabla_a S
      + \lambda_4\,n^a \nabla_a \Phi
      + \lambda_2 v_\perp
     \Bigr).
\end{align}
For the spatial components $v^i$ one obtains
\begin{equation}
\Pi_i
  := \frac{\partial\mathcal{L}_{\mathrm{RSVP}}}{\partial(\partial_t v^i)}
   = \sqrt{h}\,\beta_v\,\omega_{i\perp}
   + \sqrt{h}\,(\text{terms from couplings}).
\end{equation}
The temporal component $v_\perp$ appears algebraically and yields a constraint rather than an independent momentum.

\subsection{Hamiltonian density}

The Hamiltonian density is given by
\[
\mathcal{H}_{\mathrm{RSVP}}
  = \Pi_\Phi \partial_t \Phi
    + \Pi_S \partial_t S
    + \Pi_i \partial_t v^i
    - \mathcal{L}_{\mathrm{RSVP}},
\]
which, after standard manipulations, can be written in ADM form
\begin{equation}
\mathcal{H}_{\mathrm{RSVP}}
  = N \mathcal{H}_\perp + N^i \mathcal{H}_i,
\end{equation}
where $\mathcal{H}_\perp$ is the Hamiltonian constraint and $\mathcal{H}_i$ are the momentum constraints.

The Hamiltonian constraint contains:
\begin{itemize}
  \item the usual geometric term
    \[
      \mathcal{H}_\perp^{\mathrm{geom}}
         = \frac{1}{2\kappa}\Bigl(
              K_{ij}K^{ij} - K^2 + {}^{(3)}R
            \Bigr)
           + \Lambda_{\mathrm{eff}},
    \]
    with $K_{ij}$ the extrinsic curvature of $\Sigma_t$ and ${}^{(3)}R$ its scalar curvature,
  \item plus RSVP contributions
    \[
    \mathcal{H}_\perp^{\mathrm{RSVP}}
      = \frac{1}{2\sqrt{h}} \Bigl(
          \Pi_\Phi^2/\alpha_\Phi
        + \Pi_S^2/\alpha_S
        + \Pi_i \Pi^i/\beta_v
        + \text{cross terms from }\lambda_i
        \Bigr) + \sqrt{h}\,\mathcal{V}_{\mathrm{RSVP}},
    \]
    where $\mathcal{V}_{\mathrm{RSVP}}$ collects spatial gradient terms of $\Phi,S,v^i$ and the potential $V(\Phi,S)$.
\end{itemize}

The momentum constraint $\mathcal{H}_i$ enforces spatial diffeomorphism invariance, and includes contributions from the flow of $v^i$ and the transport of $\Phi,S$ along the shift.

\subsection{Hamiltonian evolution}

Hamilton’s equations take the usual form:
\begin{align}
\partial_t \Phi &= \frac{\delta \mathcal{H}}{\delta \Pi_\Phi}, &
\partial_t \Pi_\Phi &= -\frac{\delta \mathcal{H}}{\delta \Phi}, \\
\partial_t S &= \frac{\delta \mathcal{H}}{\delta \Pi_S}, &
\partial_t \Pi_S &= -\frac{\delta \mathcal{H}}{\delta S}, \\
\partial_t v^i &= \frac{\delta \mathcal{H}}{\delta \Pi_i}, &
\partial_t \Pi_i &= -\frac{\delta \mathcal{H}}{\delta v^i},
\end{align}
subject to the constraint surfaces $\mathcal{H}_\perp=0$ and $\mathcal{H}_i=0$.

The resulting Hamiltonian system defines the RSVP plenum as a constrained Lorentzian field theory, with the entropy flow encoded in the sign and structure of the couplings $\lambda_i$ and the positivity of $\sigma_{\mathrm{RSVP}}$.

\section{Lorentzian Signature and Constructive Logic}

\subsection{Signature choice}

The choice of Lorentzian signature $(-,+,+,+)$ is essential for encoding causal structure:
\begin{itemize}
  \item timelike directions define physical evolution and the entropic arrow,
  \item null cones delineate the domain of influence and support the horizon–thermodynamic analogy,
  \item the RSVP flow $v^a$ can be taken timelike or null, depending on the regime.
\end{itemize}

The sign convention in \eqref{eq:L-fields} ensures that the energy density extracted from $\mathcal{H}_\perp$ is bounded below when the entropy production $\sigma_{\mathrm{RSVP}}$ is non–negative and the potential $V$ has a stable minimum at the RSVP ground state.

A Euclideanised version obtained by Wick rotation $t\mapsto -i\tau$ would change the sign of the kinetic terms and is useful for certain path–integral or renormalisation–group analyses, but the physical interpretation of entropy flow and causal structure rests on the Lorentzian formulation.

\subsection{Constructive interpretation in homotopy type theory}

The RSVP plenum, as a Lorentzian field theory, can be internalised in a homotopy type theory with a cohesive or synthetic differential geometric layer. The basic elements of the constructive interpretation are:
\begin{itemize}
  \item events as points of a type $\mathsf{M}$,
  \item causal curves as paths in $\mathsf{M}$ with additional structure expressing timelike or null character,
  \item fields $(\Phi,v,S)$ as dependent types (or sections of bundles) over $\mathsf{M}$,
  \item the action functional $\mathcal{S}_{\mathrm{RSVP}}$ as a functional type assigning a real value to each triple $(\Phi,v,S)$.
\end{itemize}

The Euler–Lagrange equations are represented not as classical equalities on real–valued functions but as inhabited types of proofs
\[
\mathsf{EOM}_\Phi(\Phi,v,S,g),\quad
\mathsf{EOM}_v(\Phi,v,S,g),\quad
\mathsf{EOM}_S(\Phi,v,S,g),\quad
\mathsf{EOM}_g(\Phi,v,S,g),
\]
each asserting stationarity of the action with respect to the corresponding variation. Solutions are proofs–as–programs witnessing that a given configuration satisfies these types.

There is no conflict between Lorentzian signature and constructive logic:
\begin{itemize}
  \item the causal structure is encoded as a relation on $\mathsf{M}$ with proofs witnessing causal links,
  \item global hyperbolicity, existence of Cauchy surfaces, and other structural properties are expressed as higher–inductive or modal conditions on types, rather than as classical existence theorems relying on excluded middle,
  \item univalence identifies isomorphic Lorentzian plenum models (e.g.\ different foliations, gauge–equivalent fields) as equal at the level of types, reflecting the physical irrelevance of redundant structure.
\end{itemize}

In particular, the Hamiltonian constraint $\mathcal{H}_\perp=0$ becomes a type $\mathsf{HamConstraint}$ whose inhabitants are certificates that a given slice $(h_{ij},K_{ij},\Phi,S,v^i,\ldots)$ is admissible. The constructive semantics require explicit construction of such certificates rather than non–constructive existence proofs.

\subsection{HoTT modal structure and causal modalities}

Modal operators in the type theory can be used to encode causal and entropic restrictions:
\begin{itemize}
  \item a ``future'' modality $\lozenge_{\mathrm{future}}$ capturing what can be reached along future–directed causal curves,
  \item a ``past'' modality $\square_{\mathrm{past}}$ expressing invariants under past light–cone restriction,
  \item the RSVP modalities $\Box_{\mathrm{RSVP}}$ and $\Box_{\mathrm{Polyx}}$ interpreted as stability under field evolution and EHR rewriting.
\end{itemize}
Compositions of these modalities express statements such as ``this property holds for all future harmonic evolutions compatible with the causal structure'', which align with the interpretation of the Grand Irreversibility Theorem in a constructive, proof–relevant setting.

\section{Summary}

The RSVP plenum is a Lorentzian field theory on $(\mathcal{M},g_{ab})$ with action
\[
\mathcal{S}_{\mathrm{RSVP}}[g,\Phi,v,S]
  = \int \sqrt{-g}\,
    \Bigl(
      \tfrac{1}{2\kappa}R
      -\Lambda_{\mathrm{eff}}
      -\tfrac{\alpha_\Phi}{2}\nabla\Phi\cdot\nabla\Phi
      -\tfrac{\alpha_S}{2}\nabla S\cdot\nabla S
      -\tfrac{\alpha_v}{2} h_{ab} v^a v^b
      -\tfrac{\beta_v}{2}\omega^2
      -V(\Phi,S)
      + \mathcal{L}_{\mathrm{coupling}}
    \Bigr)\,d^4x,
\]
whose Euler–Lagrange equations yield the RSVP dynamics and whose 3+1 decomposition defines a constrained Hamiltonian system with well–behaved energy density.

The Lorentzian signature provides the causal and entropic structure of the plenum and is fully compatible with a constructive, homotopy–type–theoretic interpretation in which fields, actions, and solutions are proof–relevant objects. The constructive layer enforces that all claims about the plenum be backed by explicit witnesses (proofs), while univalence identifies physically equivalent configurations. Together, the Lagrangian–Hamiltonian formalism and the HoTT interpretation present RSVP as a Lorentzian, causally ordered, and constructively verified theory of the plenum.

